% GENERATED FILE. Edit canonical lessons, then run npm run generate:books.
% canonical-source-hash: fnv1a64:4e8be60218993b36
% canonical-lessons: GU-C202-corvu, GU-C202-thanu, GU-C202-luntvu, GU-C202-jadvu, GU-C202-kadhvu

\chapter{A theft in the bazaar}
\label{ch:gu-a-theft-in-the-bazaar}
\hlchaptermodality{202}

\begin{chapteropening}
\textbf{By the end of this chapter:} \emph{I can report a theft at the police station.}

Complete the last lesson of chapter 202: I can report a theft at the police station.
\end{chapteropening}

\section[\texorpdfstring{\gu{ચોરવું} (corvũ)}{ચોરવું (corvũ)}]{\gu{ચોરવું} (corvũ) — to steal}

\label{lesson:GU-C202-corvu}



\begin{quote}
Before the new one: say the Gujarati for barber, then the Gujarati for carpenter.
\end{quote}

\subsection*{You'll want to know: \gu{ચોરવું}}
\textbf{\gu{ચોરવું}} — \emph{corvũ} — \textquotedblleft{}to steal\textquotedblright{}.

\begin{quote}
\textbf{\gu{કોઈએ} \gu{મારું} \gu{પાકીટ} \gu{ચોર્યું}.} — \emph{koīe mārũ pākīṭ coryũ} — \textquotedblleft{}Someone stole my wallet.\textquotedblright{}
\end{quote}

\begin{cousinweb}
\textbf{\gu{ચોરવું}} (\emph{corvũ}) is a tadbhava. It comes from Sanskrit \emph{cōráyati}, \textquotedblleft{}steals\textquotedblright{}, and CDIAL glosses the Gujarati verb \textquotedblleft{}to steal\textquotedblright{} (CDIAL 4933).
\end{cousinweb}

\begin{grammarlens}[title={What you've built}]
One more word for everyday life at A2.
\end{grammarlens}

\subsection*{Your turn}
\begin{itemize}
\raggedright
  \item \emph{Say it:} \emph{corvũ}
  \item \emph{Say it:} \emph{corvũ}, once more
  \item \emph{Say it:} \emph{koīe mārũ pākīṭ coryũ} — Someone stole my wallet.
  \item \emph{Recall:} say the Gujarati for barber, then the Gujarati for carpenter, then say \emph{corvũ} again
\end{itemize}

\begin{tcolorbox}[breakable,colback=teal!4,colframe=teal!35!black,title={Before you move on}]
What does \gu{ચોરવું} mean? (\textquotedblleft{}To steal\textquotedblright{}.) Say it once more.
\end{tcolorbox}
\section[\texorpdfstring{\gu{થાણું} (thāṇũ)}{થાણું (thāṇũ)}]{\gu{થાણું} (thāṇũ) — police station}

\label{lesson:GU-C202-thanu}



\begin{quote}
Before the new one: say the Gujarati for carpenter, then the Gujarati for to steal.
\end{quote}

\subsection*{You'll want to know: \gu{થાણું}}
\textbf{\gu{થાણું}} — \emph{thāṇũ} — \textquotedblleft{}police station\textquotedblright{}.

\begin{quote}
\textbf{\gu{પોલીસ} \gu{થાણું} \gu{સ્ટેશન} \gu{પાસે} \gu{છે}.} — \emph{polīs thāṇũ sṭeśan pāse che} — \textquotedblleft{}The police station is near the railway station.\textquotedblright{}
\end{quote}

\begin{cousinweb}
\textbf{\gu{થાણું}} (\emph{thāṇũ}) is a tadbhava. CDIAL lists it, glossed \textquotedblleft{}place, police station\textquotedblright{}, under Sanskrit \emph{sth\'{\={a}}na}, \textquotedblleft{}firm stance; place of abode\textquotedblright{} (CDIAL 13753).
\end{cousinweb}

\begin{grammarlens}[title={What you've built}]
One more word for everyday life at A2.
\end{grammarlens}

\subsection*{Your turn}
\begin{itemize}
\raggedright
  \item \emph{Say it:} \emph{thāṇũ}
  \item \emph{Say it:} \emph{thāṇũ}, once more
  \item \emph{Say it:} \emph{polīs thāṇũ sṭeśan pāse che} — The police station is near the railway station.
  \item \emph{Recall:} say the Gujarati for carpenter, then the Gujarati for to steal, then say \emph{thāṇũ} again
\end{itemize}

\begin{tcolorbox}[breakable,colback=teal!4,colframe=teal!35!black,title={Before you move on}]
What does \gu{થાણું} mean? (\textquotedblleft{}Police station\textquotedblright{}.) Say it once more.
\end{tcolorbox}
\section[\texorpdfstring{\gu{લૂંટવું} (lū̃ṭvũ)}{લૂંટવું (lū̃ṭvũ)}]{\gu{લૂંટવું} (lū̃ṭvũ) — to rob}

\label{lesson:GU-C202-luntvu}



\begin{quote}
Before the new one: say the Gujarati for to steal, then the Gujarati for police station.
\end{quote}

\subsection*{You'll want to know: \gu{લૂંટવું}}
\textbf{\gu{લૂંટવું}} — \emph{lū̃ṭvũ} — \textquotedblleft{}to rob\textquotedblright{}.

\begin{quote}
\textbf{\gu{કોઈએ} \gu{વેપારીને} \gu{લૂંટ્યો}.} — \emph{koīe vepārīne lū̃ṭyo} — \textquotedblleft{}Someone robbed the merchant.\textquotedblright{}
\end{quote}

\begin{cousinweb}
\textbf{\gu{લૂંટવું}} (\emph{lū̃ṭvũ}) is a tadbhava. CDIAL lists it, glossed \textquotedblleft{}robs\textquotedblright{}, under \emph{*luṭṭati}, \textquotedblleft{}plunders\textquotedblright{}, a reconstructed form (CDIAL 11078).
\end{cousinweb}

\begin{grammarlens}[title={What you've built}]
One more word for everyday life at A2.
\end{grammarlens}

\subsection*{Your turn}
\begin{itemize}
\raggedright
  \item \emph{Say it:} \emph{lū̃ṭvũ}
  \item \emph{Say it:} \emph{lū̃ṭvũ}, once more
  \item \emph{Say it:} \emph{koīe vepārīne lū̃ṭyo} — Someone robbed the merchant.
  \item \emph{Recall:} say the Gujarati for to steal, then the Gujarati for police station, then say \emph{lū̃ṭvũ} again
\end{itemize}

\begin{tcolorbox}[breakable,colback=teal!4,colframe=teal!35!black,title={Before you move on}]
What does \gu{લૂંટવું} mean? (\textquotedblleft{}To rob\textquotedblright{}.) Say it once more.
\end{tcolorbox}
\section[\texorpdfstring{\gu{જડવું} (jaḍvũ)}{જડવું (jaḍvũ)}]{\gu{જડવું} (jaḍvũ) — to be found, to turn up}

\label{lesson:GU-C202-jadvu}



\begin{quote}
Before the new one: say the Gujarati for police station, then the Gujarati for to rob.
\end{quote}

\subsection*{You'll want to know: \gu{જડવું}}
\textbf{\gu{જડવું}} — \emph{jaḍvũ} — \textquotedblleft{}to be found, to turn up\textquotedblright{}.

\begin{quote}
\textbf{\gu{મને} \gu{પાકીટ} \gu{જડ્યું}.} — \emph{mane pākīṭ jaḍyũ} — \textquotedblleft{}I found my wallet.\textquotedblright{}
\end{quote}

The finder takes \textbf{-\gu{ને}}: \textbf{\gu{મને} \gu{જડ્યું}}.

\begin{cousinweb}
\textbf{\gu{જડવું}} (\emph{jaḍvũ}) is listed by CDIAL as \emph{jaṛvũ}, \textquotedblleft{}to join, meet with, set jewels\textquotedblright{}, under \emph{*jaḍati}, \textquotedblleft{}joins, sets\textquotedblright{}, a word that CDIAL takes from Dravidian, citing Bloch (CDIAL 5091). Of its senses, \textquotedblleft{}meet with\textquotedblright{} is the nearest to \textquotedblleft{}to be found\textquotedblright{}.
\end{cousinweb}

\begin{grammarlens}[title={What you've built}]
One more word for everyday life at A2.
\end{grammarlens}

\subsection*{Your turn}
\begin{itemize}
\raggedright
  \item \emph{Say it:} \emph{jaḍvũ}
  \item \emph{Say it:} \emph{jaḍvũ}, once more
  \item \emph{Say it:} \emph{mane pākīṭ jaḍyũ} — I found my wallet.
  \item \emph{Recall:} say the Gujarati for police station, then the Gujarati for to rob, then say \emph{jaḍvũ} again
\end{itemize}

\begin{tcolorbox}[breakable,colback=teal!4,colframe=teal!35!black,title={Before you move on}]
What does \gu{જડવું} mean? (\textquotedblleft{}To be found, to turn up\textquotedblright{}.) Say it once more.
\end{tcolorbox}
\section[\texorpdfstring{\gu{કાઢવું} (kāḍhvũ)}{કાઢવું (kāḍhvũ)}]{\gu{કાઢવું} (kāḍhvũ) — to take out}

\label{lesson:GU-C202-kadhvu}



\begin{quote}
Before the new one: say the Gujarati for to rob, then the Gujarati for to be found.
\end{quote}

\subsection*{You'll want to know: \gu{કાઢવું}}
\textbf{\gu{કાઢવું}} — \emph{kāḍhvũ} — \textquotedblleft{}to take out\textquotedblright{}.

\begin{quote}
\textbf{\gu{તેણે} \gu{થેલીમાંથી} \gu{પૈસા} \gu{કાઢ્યા}.} — \emph{teṇe thelīmā̃thī paisā kāḍhyā} — \textquotedblleft{}He took the money out of the bag.\textquotedblright{}
\end{quote}

\begin{cousinweb}
\textbf{\gu{કાઢવું}} (\emph{kāḍhvũ}) is a tadbhava. CDIAL lists it, glossed \textquotedblleft{}to pull out\textquotedblright{}, under \emph{*kaḍḍhati}, \textquotedblleft{}pulls, draws\textquotedblright{}, a reconstructed form (CDIAL 2660).
\end{cousinweb}

\begin{grammarlens}[title={What you've built}]
One more word for everyday life at A2.
\end{grammarlens}

\subsection*{Your turn}
\begin{itemize}
\raggedright
  \item \emph{Say it:} \emph{kāḍhvũ}
  \item \emph{Say it:} \emph{kāḍhvũ}, once more
  \item \emph{Say it:} \emph{teṇe thelīmā̃thī paisā kāḍhyā} — He took the money out of the bag.
  \item \emph{Recall:} say the Gujarati for to rob, then the Gujarati for to be found, then say \emph{kāḍhvũ} again
\end{itemize}

\begin{tcolorbox}[breakable,colback=teal!4,colframe=teal!35!black,title={Before you move on}]
What does \gu{કાઢવું} mean? (\textquotedblleft{}To take out\textquotedblright{}.) Say it once more.
\end{tcolorbox}
